\textbf{Lösung zu Übung 5, Aufgabe 2}

\begin{loesung}
\begin{enumerate}[label=\alph*)]
\item
\[
G = (\,E \mid A\,) =
\begin{pmatrix}
1 & & & & 1\\
& 1 & & & 1\\
& & 1 & & 1\\
& & & 1 & 1
\end{pmatrix}, \qquad
H = (-A^T \mid E) = (\,2\;2\;2\;2 \mid 1\,)
\]
% KORRIGIERT: "a G = 2 1 2 0 G" (a nicht fett, Klammern fehlten)
\item $\mathbf{a}\,G = (\,2\;1\;2\;0\,)\,G = 2\;1\;2\;0\;2$ \quad (Paritätsstelle $2+1+2+0 = 5 \equiv 2 \pmod{3}$)
\item
% KORRIGIERT: Syndrom ist hier ein Skalar in F_3; "\ne \mathbf{o}" -> "\ne 0" (einheitlich mit der letzten Zeile)
\begin{align*}
S(\mathbf{y}_1) &= H\,\mathbf{y}_1^T = (\,2\;2\;2\;2\;\;1\,)\;\mathbf{y}_1^T = 2 \ne 0\\
S(\mathbf{y}_2) &= H\,\mathbf{y}_2^T = (\,2\;2\;2\;2\;\;1\,)\;\mathbf{y}_2^T = 1 \ne 0\\
S(\mathbf{y}_3) &= H\,\mathbf{y}_3^T = (\,2\;2\;2\;2\;\;1\,)\;\mathbf{y}_3^T = 0 \;\Rightarrow\; \mathbf{y}_3 \in C
\end{align*}
\end{enumerate}
\end{loesung}
